% AXIOM TEMPLATE: ieee-journal
% !! AXIOM-LOCKED-PREAMBLE-START !!  —— preamble 由模板锁定，agent 不得修改。
% 禁止改 \documentclass、禁止删除已有 \usepackage/\usetikzlibrary。
% Tectonic 会自动从 CTAN 下载 IEEEtran.cls (首次编译联网)。
\documentclass[journal]{IEEEtran}

% —— 数学 ——
\usepackage{amsmath,amssymb,amsthm,mathtools}
\usepackage{bm}

% —— 图形与表格 ——
\usepackage{graphicx}
\usepackage{xcolor}
\usepackage{booktabs}
\usepackage{array}
\usepackage{multirow}
\usepackage{tabularx}

% —— TikZ (常用库全加载) ——
\usepackage{tikz}
\usetikzlibrary{positioning,arrows.meta,shapes.geometric,shapes.misc,calc,fit,backgrounds,decorations.pathreplacing}

% —— 伪代码 ——
\usepackage{algorithm}
\usepackage{algpseudocode}

% —— 列表 ——
\usepackage{enumitem}

% —— 引用与链接 ——
\usepackage{url}
\usepackage{hyperref}
\hypersetup{colorlinks=true,linkcolor=blue,citecolor=blue,urlcolor=blue}

% —— 微排版 ——
\usepackage{microtype}
% !! AXIOM-LOCKED-PREAMBLE-END !!

% AXIOM_BODY: agent 只能从 \begin{document} 之后开始编辑正文。
\title{Paper Title}

\author{
  Author~Name,~\IEEEmembership{Member,~IEEE}
  \thanks{Manuscript received...; revised... (affiliation, email).}%
}

\begin{document}
\maketitle

\begin{abstract}
Replace this with the abstract.
\end{abstract}

\begin{IEEEkeywords}
keyword1, keyword2, keyword3
\end{IEEEkeywords}

\section{Introduction}
Replace this with the introduction.

% 引用: 写好 references.bib 后, 正文用 \cite{key}, 并取消下面两行注释。
% 未引用任何文献时保持注释, 否则空 .bbl 会触发 "missing \item" 报错。
% \bibliographystyle{IEEEtran}
% \bibliography{references}
\end{document}
